\section{Dataset Used in \benchmarkname}\label{sec:dataset_appendix}

\begin{table}[ht]
\centering
\caption{Statistics of the knowledge graphs used for evaluation}
\label{tab:datasets}
% \resizebox{0.9\columnwidth}{!}{%
\begin{tabular}{@{}cccc@{}}
\toprule
\textbf{Attributes}     & \textbf{Academia} & \textbf{Literature} & \textbf{E-commerce} \\ \midrule
%\textbf{Domain}         & Academia         & Literature         & E-commerce      \\
% \textbf{Entity Types}   & 3                & 4                  & 2               \\
% \textbf{Relation Types} & 6                & 7                  & 5               \\
\# of Entities       & 2.7M             & 3.7M               & 8.1M            \\
\# of Relations      & 67M              & 28M                & 193M            \\
\# Templates      & 24              & 28                & 21  \\
\# of Questions      & 400              & 400                & 400  \\ \bottomrule
\end{tabular}%
% }
\end{table}

The statistics of the datasets used in \benchmarkname{} are presented in~\autoref{tab:datasets}. The three knowledge graphs are sourced from RGBench~\cite{graphcot}. Specifically, in the academic domain, there are three types of entities: \textit{papers}, \textit{authors}, and \textit{venues}. These are naturally interconnected by citation-related relationships such as ``\textit{written-by}'' and ``\textit{publish-in}''. The selected academic graph is from the Physics domain of DBLP~\cite{academia}. 

In the literature domain, the knowledge graph contains four types of entities—\textit{books}, \textit{authors}, \textit{publishers}, and \textit{series}. Its structure captures the publication relationships among these entities, extracted from the Goodreads dataset~\cite{literature}, which provides a rich collection of books along with detailed metadata. The relations include ``\textit{written-by}'', ``\textit{publish-in}'', ``\textit{book-series}'', and others.

For the e-commerce domain, we adopt the graph from the Amazon product dataset~\cite{e-commerce}, which offers metadata for items across a wide range of product categories. In this domain, the nodes include \textit{items} and \textit{brands}, and the interconnecting relations include ``\textit{also-viewed}``, ``\textit{also-bought}'', ``\textit{buy-after-viewing}'', ``\textit{bought-together}'' and ``\textit{item-brand}''.


In the following, we list the concrete graph schema (i.e., entity and relation types) in the above three knowledge graphs.

\subsection{Academia Domain}

\textit{Node properties:}
\squishlist
\item type: author, properties: ["id", "name", "node\_type"]
\item type: paper, properties: ["id", "name", "label", "year", "node\_type", "abstract"]
\item type: venue, properties: ["id", "name", "node\_type"]
\squishend

\textit{Edge properties:}

Author nodes are linked to their paper nodes by authorship.
\squishlist
\item author $\rightarrow$ "paper" $\rightarrow$ paper
\squishend

Paper nodes are linked to their author nodes, venue nodes, other paper nodes.
\squishlist
\item paper $\rightarrow$ "author" $\rightarrow$ author
\item paper $\rightarrow$ "reference" $\rightarrow$ paper
\item paper $\rightarrow$ "venue" $\rightarrow$ venue
\item paper $\rightarrow$ "cited\_by" $\rightarrow$ paper
\squishend

Venue nodes are linked to their included paper nodes.
\squishlist
\item venue $\rightarrow$ "paper" $\rightarrow$ paper
\squishend

\subsection{Literature Domain}

\textit{Node properties:}
\squishlist
\item type: book, properties: ["id", "name", "node\_type", "description", "publication\_year", "genres"]
\item type: author, properties: ["id", "name", "node\_type"]
\item type: publisher, properties: ["id", "name", "node\_type"]
\item type: series, properties: ["id", "name", "node\_type", "description"]
\squishend

\textit{Edge properties:}

Book nodes are linked to book nodes, author nodes, publisher nodes and series nodes.
\squishlist
\item book $\rightarrow$ "author" $\rightarrow$ author
\item book $\rightarrow$ "publisher" $\rightarrow$ publisher
\item book $\rightarrow$ "series" $\rightarrow$ series
\item book $\rightarrow$ "similar\_books" $\rightarrow$ book
\squishend

Author nodes are linked to their neighboring book nodes.
\squishlist
\item author $\rightarrow$ "book" $\rightarrow$ book
\squishend

Publisher nodes are linked to their neighboring book nodes.
\squishlist
\item publisher $\rightarrow$ "book" $\rightarrow$ book
\squishend

Series nodes are linked to their neighboring book nodes.
1. series $\rightarrow$ "book" $\rightarrow$ book

\subsection{E-commerce Domain}

\textit{Node properties:}
\squishlist
\item type: item, properties: ["id", "name", "node\_type"]
\item type: brand, properties: ["id", "name", "node\_type"]
\squishend

Edge properties:
Item nodes are linked to neighboring item nodes and brand nodes.
\squishlist
\item item $\rightarrow$ "also\_viewed\_item" $\rightarrow$ item
\item item $\rightarrow$ "buy\_after\_viewing\_item" $\rightarrow$ item
\item item $\rightarrow$ "also\_bought\_item" $\rightarrow$ item
\item item $\rightarrow$ "bought\_together\_item" $\rightarrow$ item
\item item $\rightarrow$ "brand" $\rightarrow$ brand
\squishend

Brand nodes are linked to their neighboring item nodes. Specific relations are:
\squishlist
\item brand $\rightarrow$ "item" $\rightarrow$ item
\squishend